\documentclass[A4paper,11pt]{report}

% ================== PACKAGES ==================
\usepackage{lmodern}
\usepackage[T1]{fontenc}
\usepackage[french]{babel}

\usepackage[Conny]{fncychap}
\usepackage{titlesec}

\usepackage{fancybox}
\usepackage[top=2cm, bottom=2cm, left=2cm, right=2cm]{geometry}
\usepackage{multicol}
\setlength{\columnsep}{1cm}

\usepackage{amsmath,amsfonts,amsthm,amssymb,mathtools}
\usepackage{enumitem}
\usepackage{framed}
\usepackage{graphicx}

\usepackage[section]{placeins}

\makeatletter
\setlength{\@fptop}{0pt}
\makeatother

\usepackage{fancyhdr}
\usepackage{tikz}
\usepackage{circuitikz}

\usepackage{fourier-orns}
\usepackage{lastpage}
\usepackage{mwe}
\usepackage{afterpage}

\usepackage{hyperref} % toujours en dernier

\usepackage{listings}


% --- Couleurs Ubuntu ---
\definecolor{ubuntubg}{RGB}{48,10,36}
\definecolor{ubuntutext}{RGB}{238,238,236}
\definecolor{ubuntuorange}{RGB}{233,84,32}
\definecolor{ubuntugreen}{RGB}{78,154,6}
\definecolor{ubuntugray}{RGB}{136,138,133}

% --- Couleurs VSCode (dark+) ---
\definecolor{vscodebg}{RGB}{30,30,30}
\definecolor{vscodetext}{RGB}{212,212,212}
\definecolor{vscodeblue}{RGB}{86,156,214}
\definecolor{vscodegreen}{RGB}{106,153,85}
\definecolor{vscodeorange}{RGB}{206,145,120}
\definecolor{vscodepurple}{RGB}{197,134,192}
\definecolor{vscodegray}{RGB}{128,128,128}

\lstdefinestyle{ubuntu}{
    backgroundcolor=\color{ubuntubg},   
    basicstyle=\ttfamily\footnotesize\color{ubuntutext},
    keywordstyle=\color{ubuntuorange},
    commentstyle=\color{ubuntugreen},
    stringstyle=\color{ubuntugray},
    numbers=none,
    breaklines=true,
    breakatwhitespace=false,
    keepspaces=true,
    showspaces=false,
    showstringspaces=false,
    showtabs=false,
    frame=single,
    rulecolor=\color{ubuntugray},
    tabsize=2
}

\lstdefinestyle{vscodeC}{
    language=C,
    backgroundcolor=\color{vscodebg},
    basicstyle=\ttfamily\footnotesize\color{vscodetext},
    keywordstyle=\color{vscodeblue}\bfseries,
    commentstyle=\color{vscodegreen}\itshape,
    stringstyle=\color{vscodeorange},
    numberstyle=\tiny\color{vscodegray},
    numbers=left,
    stepnumber=1,
    numbersep=10pt,
    breaklines=true,
    breakatwhitespace=false,
    keepspaces=true,
    showspaces=false,
    showstringspaces=false,
    showtabs=false,
    frame=none,
    rulecolor=\color{vscodegray},
    tabsize=4
}

\lstset{style=vscodeC}

% ================== TIKZ STYLE ==================
\tikzset{every node/.style={font=\large}}

% ================== NUMEROTATION ==================
\renewcommand{\thesection}{\Roman{section}.}
\renewcommand{\thesubsection}{Q\arabic{subsection}.}
\renewcommand{\thesubsubsection}{\arabic{subsubsection}.}

\setcounter{secnumdepth}{3}

% ================== FOOTER DECO ==================
\renewcommand\footrule{%
\hrulefill}

% ================== STYLE PLAIN ==================
\fancypagestyle{plain}{%
  \fancyhf{}
  \fancyfoot[C]{\thepage\ of \pageref{LastPage}}
  \renewcommand{\headrulewidth}{0pt}
  \renewcommand{\footrulewidth}{0.4pt}
}

% ================== STYLE FANCY ==================
\fancypagestyle{fancy}{%
  \fancyhf{}
  \fancyfoot[C]{\thepage\ / \pageref{LastPage}}
  \fancyhead[L]{BRIENT Tidiane \\ NYTUN Jesper}
  \fancyhead[C]{\textit{ TP C et réseaux }}
  \fancyhead[R]{3 IMACS D}
  \renewcommand{\headrulewidth}{0.4pt}
  \renewcommand{\footrulewidth}{0.4pt}
}

\pagestyle{fancy}
\setlength{\fboxsep}{0.5cm}




\newcommand{\itemb}{\item[$\bullet$ ] }
% ================== DOCUMENT ==================
\begin{document}

\section*{\huge \centering Compte-Rendu de TP Langage C et Réseaux}

\vspace{2cm}
\textbf{Intro rapide:} Ce TP met en jeu l'implémentation de la fonction tsock, permettant l'envoi de messages via internet. 

\section{Implémentation de tsock-v1}

\subsubsection{Prise en compte de l'option -u}

Une fois le fichier \textit{tsock-v0.c} dupliqué et rennomé en \textit{tsock-v1.c} , on commence par ajouter une variable \textit{protocole} initialisée à 0 qui nous indiquera le protocole utilisé:
\begin{itemize}
    \itemb 1 si le protocole utilisé est UDP (-u présent)
    \itemb 0 si TCP est utilisé (option par défaut, voir partie 2 du travail)
\end{itemize}

Pour cela, on rajoute dans la première boucle while:

\begin{lstlisting}
    case 'u':
			protocole =1;		
			break;     

\end{lstlisting}

de manière à affecter la valeur 1 à protocole lorsque UDP est utilisé.\\


Ensuite, on ajoute un if à la fin du programme permettant, si un $-u$ est présent, d'éxécuter les actions relatives au protocole UDP dans le cas d'une source ou dans le cas d'un puit, soit:
\begin{itemize}
    \itemb Dans le cas d'une source:
    \begin{itemize}
        \item créer un socket UDP;
        \item construire l'adresse du socket cible;
        \item envoi des données.
    \end{itemize}
    \itemb Dans le cas d'un puits:
    \begin{itemize}
        \item créer un socket UDP;
        \item assigner une adresse au socket crée;
        \item se mettre en écoute.
    \end{itemize}
\end{itemize} 



Ce qui se traduit en code par:

\begin{lstlisting}
    if(protocole ==1){
		if(source==1){
			int sock = socket(AF_INET,SOCK_DGRAM,IPPROTO_UDP);
            /*Ajouter le code*/
		}	
        else{
			int sock = socket(AF_INET,SOCK_DGRAM,IPPROTO_UDP);
        /*Ajouter le code*/
		}
	}
\end{lstlisting}



\end{document}
